Bagian ini menjelaskan struktur source program yang digunakan untuk menyelesaikan permasalahan \textit{Ice Sliding Puzzle Solver}. 
Program dikembangkan menggunakan bahasa C++ dengan pendekatan modular, sehingga komponen pembacaan input, representasi peta, simulasi pergerakan, algoritma pathfinding, heuristik, CLI, dan GUI dipisahkan ke dalam folder dan class yang berbeda.

\begin{tcolorbox}[
	colback=blue!5,
	colframe=blue!60!black,
	title=\textbf{Ringkasan Implementasi Program},
	arc=2mm,
	boxrule=0.8pt
]
Program membaca konfigurasi papan dari file \texttt{.txt}, memvalidasi isi papan, membangun representasi map dan cost, lalu menjalankan algoritma pathfinding yang dipilih pengguna. 
Algoritma yang diimplementasikan adalah \textbf{Uniform Cost Search}, \textbf{Greedy Best-First Search}, \textbf{A*}, dan \textbf{Weighted A*}. 
Hasil akhir program berupa solusi gerakan, total cost, jumlah iterasi, waktu eksekusi, visualisasi langkah solusi, serta opsi penyimpanan solusi ke file.
\end{tcolorbox}

\section{Informasi Umum Program}

\begin{table}[H]
	\centering
	\caption{Informasi umum program}
	\begin{tabular}{|p{4cm}|p{10cm}|}
		\hline
		\rowcolor{blue!15}
		\textbf{Aspek} & \textbf{Keterangan} \\
		\hline
		Bahasa pemrograman & C++ \\
		\hline
		Standar bahasa & C++17 \\
		\hline
		Alasan pemilihan bahasa & C++ dipilih karena mendukung pemrograman berorientasi objek, memiliki performa baik untuk eksplorasi state space, serta menyediakan struktur data standar seperti \texttt{priority\_queue}, \texttt{map}, \texttt{vector}, dan \texttt{tuple}. \\
		\hline
		Compiler & \texttt{g++} dengan standar kompilasi \texttt{-std=c++17}. \\
		\hline
		Library standar & \texttt{iostream}, \texttt{fstream}, \texttt{vector}, \texttt{string}, \texttt{queue}, \texttt{map}, \texttt{tuple}, \texttt{chrono}, dan library standar C++ lainnya. \\
		\hline
		Library eksternal & Raylib digunakan apabila program dijalankan dengan fitur GUI. Jika Raylib tidak tersedia, program tetap dapat dijalankan melalui CLI. \\
		\hline
		Cara menjalankan umum & Program dikompilasi menggunakan \texttt{make}, kemudian dijalankan melalui mode CLI atau GUI sesuai target yang tersedia pada \texttt{Makefile}. \\
		\hline
	\end{tabular}
\end{table}

\section{Struktur Folder Repository}

Struktur repository dibuat modular agar setiap bagian program memiliki tanggung jawab yang jelas. 
Folder \texttt{include/} digunakan untuk menyimpan header file, sedangkan folder \texttt{src/} digunakan untuk menyimpan implementasi setiap modul.

\begin{tcolorbox}[
	colback=gray!5,
	colframe=black!60,
	title=\textbf{Struktur Folder Utama},
	arc=2mm,
	boxrule=0.8pt
]
\scriptsize
\begin{verbatim}
STIMA-TUCIL3/
|-- assets/
|-- bin/
|-- build/
|-- docs/
|-- include/
|   |-- cli/
|   |-- exception/
|   |-- gui/
|   |-- heuristic/
|   |   |-- Advanced.hpp
|   |   |-- Custom.hpp
|   |   |-- IHeuristic.hpp
|   |   `-- Manhattan.hpp
|   |-- model/
|   |   |-- GameMap.hpp
|   |   |-- MoveResult.hpp
|   |   |-- Position.hpp
|   |   |-- SearchResult.hpp
|   |   `-- State.hpp
|   |-- solver/
|   |   |-- AStar.hpp
|   |   |-- GBFS.hpp
|   |   |-- ISolver.hpp
|   |   |-- UCS.hpp
|   |   `-- WeightedAStar.hpp
|   `-- utils/
|-- saves/
|-- src/
|   |-- cli/
|   |-- exception/
|   |-- gui/
|   |-- heuristic/
|   |   |-- Advanced.cpp
|   |   |-- Custom.cpp
|   |   `-- Manhattan.cpp
|   |-- model/
|   |   |-- GameMap.cpp
|   |   |-- MoveResult.cpp
|   |   |-- Position.cpp
|   |   |-- SearchResult.cpp
|   |   `-- State.cpp
|   |-- solver/
|   |   |-- AStar.cpp
|   |   |-- GBFS.cpp
|   |   |-- UCS.cpp
|   |   `-- WeightedAStar.cpp
|   |-- utils/
|   `-- main.cpp
|-- test/
|-- .gitignore
|-- CMakeLists.txt
`-- README.md
\end{verbatim}
\end{tcolorbox}

\begin{longtable}{|p{3.5cm}|p{10.5cm}|}
	\caption{Peran folder pada repository} \\
	\hline
	\rowcolor{gray!20}
	\textbf{Folder / File} & \textbf{Peran} \\
	\hline
	\endfirsthead
	
	\hline
	\rowcolor{gray!20}
	\textbf{Folder / File} & \textbf{Peran} \\
	\hline
	\endhead
	
	\texttt{src/} & Berisi implementasi source code utama program, seperti parser, model, solver, heuristic, CLI, GUI, exception, dan utility. \\
	\hline
	\texttt{include/} & Berisi deklarasi class dan interface dalam bentuk header file. Struktur foldernya dibuat sejajar dengan \texttt{src/}. \\
	\hline
	\texttt{src/solver/} & Berisi implementasi algoritma pathfinding, yaitu \texttt{UCS.cpp}, \texttt{GBFS.cpp}, dan \texttt{AStar.cpp}. \\
	\hline
	\texttt{include/solver/} & Berisi interface solver dan header untuk masing-masing algoritma, yaitu \texttt{ISolver.hpp}, \texttt{UCS.hpp}, \texttt{GBFS.hpp}, dan \texttt{AStar.hpp}. \\
	\hline
	\texttt{src/model/} dan \texttt{include/model/} & Berisi model data utama, seperti posisi, state, map, hasil gerakan, dan hasil pencarian. \\
	\hline
	\texttt{src/heuristic/} dan \texttt{include/heuristic/} & Berisi interface dan implementasi heuristik yang digunakan oleh GBFS dan A*. \\
	\hline
	\texttt{src/utils/} dan \texttt{include/utils/} & Berisi komponen pendukung seperti movement engine, pembacaan file, penyimpanan solusi, atau utilitas lain. \\
	\hline
	\texttt{src/cli/} dan \texttt{include/cli/} & Berisi implementasi antarmuka command line untuk memilih input file, algoritma, heuristik, playback, dan penyimpanan solusi. \\
	\hline
	\texttt{src/gui/} dan \texttt{include/gui/} & Berisi implementasi GUI apabila fitur bonus GUI dijalankan. \\
	\hline
	\texttt{docs/} & Berisi dokumen laporan dalam bentuk LaTeX dan hasil PDF. \\
	\hline
	\texttt{saves/} & Berisi file hasil penyimpanan solusi atau output pathfinding. \\
	\hline
	\texttt{README.md} & Berisi deskripsi program, requirement, cara kompilasi, cara menjalankan, dan identitas pembuat. \\
	\hline
	\texttt{Makefile} & Berisi konfigurasi build agar program dapat dikompilasi dan dijalankan dengan perintah sederhana. \\
	\hline
\end{longtable}

\section{Struktur Source Code}

Source code dibagi berdasarkan peran modul. 
Pembagian ini bertujuan agar implementasi algoritma tidak bercampur dengan pembacaan input, validasi, antarmuka pengguna, maupun visualisasi.

\begin{table}[H]
	\centering
	\caption{Pembagian modul source code}
	\begin{tabular}{|p{4cm}|p{10cm}|}
		\hline
		\rowcolor{green!15}
		\textbf{Modul} & \textbf{Tanggung Jawab} \\
		\hline
		Model & Menyimpan struktur data utama seperti posisi aktor, state pencarian, peta permainan, hasil sliding, dan hasil pencarian. \\
		\hline
		Parser input & Membaca file \texttt{.txt}, mengambil ukuran papan, konfigurasi tile, dan matriks cost. \\
		\hline
		Validator input & Memastikan input valid, ukuran papan sesuai, jumlah baris benar, karakter tile valid, posisi awal dan tujuan tersedia, serta cost sesuai dimensi papan. \\
		\hline
		Movement engine & Mensimulasikan gerakan sliding ke arah atas, bawah, kiri, dan kanan. \\
		\hline
		Solver & Menyediakan interface pathfinding dan implementasi UCS, GBFS, serta A*. \\
		\hline
		Heuristic & Menghitung estimasi jarak atau cost menuju target berikutnya untuk GBFS dan A*. \\
		\hline
		Output / save & Menampilkan solusi, visualisasi langkah, dan menyimpan solusi ke file. \\
		\hline
		CLI & Menyediakan interaksi berbasis terminal. \\
		\hline
		GUI & Menyediakan visualisasi grafis apabila bonus GUI dijalankan. \\
		\hline
	\end{tabular}
\end{table}

\section{Modul dan Class Utama}

\begingroup
\footnotesize
\setlength{\tabcolsep}{4pt}
\renewcommand{\arraystretch}{1.25}

\begin{longtable}{
	|>{\raggedright\arraybackslash}p{3.1cm}
	|>{\raggedright\arraybackslash}p{4.7cm}
	|>{\raggedright\arraybackslash}p{5.4cm}|
}
	\caption{Class dan modul utama program} \\
	\hline
	\rowcolor{green!15}
	\textbf{Class / Modul} & \textbf{Lokasi Umum} & \textbf{Peran} \\
	\hline
	\endfirsthead
	
	\hline
	\rowcolor{green!15}
	\textbf{Class / Modul} & \textbf{Lokasi Umum} & \textbf{Peran} \\
	\hline
	\endhead
	
	\texttt{Position} 
	& \texttt{include/model/} dan \texttt{src/model/} 
	& Menyimpan koordinat baris dan kolom aktor atau tile tertentu. \\
	\hline
	
	\texttt{State} 
	& \texttt{include/model/} dan \texttt{src/model/} 
	& Merepresentasikan state pencarian, termasuk posisi, path, cost, heuristik, angka terakhir, dan collected mask. \\
	\hline
	
	\texttt{GameMap} 
	& \texttt{include/model/} dan \texttt{src/model/} 
	& Menyimpan konfigurasi papan, cost tile, posisi awal, posisi goal, dan informasi angka pada papan. \\
	\hline
	
	\texttt{MoveResult} 
	& \texttt{include/model/} dan \texttt{src/model/} 
	& Menyimpan hasil simulasi satu gerakan sliding, seperti posisi baru, cost gerakan, status valid, dan status game over. \\
	\hline
	
	\texttt{SearchResult} 
	& \texttt{include/model/} dan \texttt{src/model/} 
	& Menyimpan hasil akhir pencarian, seperti solusi ditemukan atau tidak, path solusi, total cost, waktu eksekusi, dan jumlah iterasi. \\
	\hline
	
	\texttt{InputParser} 
	& \texttt{include/utils/} atau \texttt{include/cli/} 
	& Membaca file testcase berekstensi \texttt{.txt}. \\
	\hline
	
	\texttt{InputValidator} 
	& \texttt{include/utils/} 
	& Memvalidasi konfigurasi papan dan matriks cost. \\
	\hline
	
	\texttt{MovementEngine} 
	& \texttt{include/utils/} dan \texttt{src/utils/} 
	& Mensimulasikan mekanisme sliding dan menentukan apakah suatu gerakan valid. \\
	\hline
	
	\texttt{ISolver} 
	& \texttt{include/solver/}\newline\texttt{ISolver.hpp} 
	& Interface umum untuk semua algoritma pathfinding. \\
	\hline
	
	\texttt{UCS} 
	& \texttt{include/solver/}\newline\texttt{UCS.hpp} dan \newline\texttt{src/solver/}\newline\texttt{UCS.cpp} 
	& Implementasi Uniform Cost Search. \\
	\hline
	
	\texttt{GBFS} 
	& \texttt{include/solver/}\newline\texttt{GBFS.hpp} dan \newline\texttt{src/solver/}\newline\texttt{GBFS.cpp} 
	& Implementasi Greedy Best-First Search. \\
	\hline
	
	\texttt{A*} 
	& \texttt{include/solver/}\newline\texttt{AStar.hpp} dan \newline\texttt{src/solver/}\newline\texttt{AStar.cpp} 
	& Implementasi A*. \\
	\hline
	
	\texttt{WeightedA*} 
	& \texttt{include/solver/}\newline\texttt{WeightedAStar.hpp} dan \newline\texttt{src/solver/}\newline\texttt{WeightedAStar.cpp} 
	& Implementasi Weighted A*. \\
	\hline
	
	\texttt{IHeuristic} 
	& \texttt{include/heuristic/} 
	& Interface umum untuk heuristik. \\
	\hline
	
	\texttt{SolutionWriter} 
	& \texttt{include/utils/} atau \texttt{src/utils/} 
	& Menyimpan solusi dan visualisasi langkah ke file. \\
	\hline
	
	\texttt{CLI} 
	& \texttt{include/cli/} dan \texttt{src/cli/} 
	& Menangani input pengguna melalui terminal. \\
	\hline
	
	\texttt{GUI} 
	& \texttt{include/gui/} dan \texttt{src/gui/} 
	& Menangani tampilan grafis dan playback visual. \\
	\hline
\end{longtable}

\endgroup

\section{Relasi Antar Modul}

Relasi antar modul dapat digambarkan sebagai alur dari input file sampai output solusi. 
Parser membaca testcase, validator memastikan testcase valid, kemudian \texttt{GameMap} dibentuk. 
Setelah itu, solver yang dipilih pengguna akan menjalankan pencarian dengan bantuan \texttt{MovementEngine} dan, untuk GBFS serta A*, bantuan modul \texttt{Heuristic}.

\begin{tcolorbox}[
	colback=gray!5,
	colframe=black!60,
	title=\textbf{Diagram Alur Source Program},
	arc=2mm,
	boxrule=0.8pt
]
\small
\begin{center}
\begin{tabular}{c}
\fbox{\parbox{0.82\textwidth}{
\centering
\textbf{1. Input}\\
File \texttt{.txt}, parser, validator
}} \\[0.25cm]
$\downarrow$ \\[-0.05cm]
\fbox{\parbox{0.82\textwidth}{
\centering
\textbf{2. Model Data}\\
\texttt{GameMap}, \texttt{State}, \texttt{MoveResult}, \texttt{SearchResult}
}} \\[0.25cm]
$\downarrow$ \\[-0.05cm]
\fbox{\parbox{0.82\textwidth}{
\centering
\textbf{3. Proses Pencarian}\\
\texttt{UCS}, \texttt{GBFS}, \texttt{A*}, \texttt{WeightedA*}, \texttt{MovementEngine}, \texttt{Heuristic}
}} \\[0.25cm]
$\downarrow$ \\[-0.05cm]
\fbox{\parbox{0.82\textwidth}{
\centering
\textbf{4. Output}\\
CLI, GUI, visualisasi langkah, playback, simpan solusi
}}
\end{tabular}
\end{center}
\end{tcolorbox}

\section{Penjelasan Peran Source Code}

\begin{longtable}{|p{4.5cm}|p{9.5cm}|}
	\caption{Peran source code terhadap alur program} \\
	\hline
	\rowcolor{blue!15}
	\textbf{Proses} & \textbf{Penjelasan Implementasi} \\
	\hline
	\endfirsthead
	
	\hline
	\rowcolor{blue!15}
	\textbf{Proses} & \textbf{Penjelasan Implementasi} \\
	\hline
	\endhead
	
	Membaca file input \texttt{.txt} & Program menerima path file dari pengguna, kemudian membaca baris pertama sebagai ukuran papan, $N$ baris berikutnya sebagai konfigurasi map, dan $N$ baris setelahnya sebagai matriks cost. \\
	\hline
	Memvalidasi konfigurasi papan & Program memeriksa kesesuaian ukuran, karakter tile, keberadaan posisi awal \texttt{Z}, keberadaan goal \texttt{O}, serta dimensi matriks cost. \\
	\hline
	Menyimpan representasi map & Konfigurasi papan dan cost disimpan dalam \texttt{GameMap}, sehingga solver dapat mengakses informasi tile, cost, posisi awal, posisi goal, dan angka. \\
	\hline
	Simulasi sliding & \texttt{MovementEngine::slide} menerima map, state saat ini, dan arah gerakan, lalu mengembalikan \texttt{MoveResult}. \\
	\hline
	Mendeteksi game over & Gerakan dianggap game over jika aktor melewati lava, keluar papan tanpa tertahan rintangan, atau melewati angka yang belum sesuai urutan. \\
	\hline
	Mendeteksi angka berurutan & Ketika sliding melewati tile angka, program memeriksa apakah angka tersebut sama dengan angka terakhir yang sudah dikumpulkan ditambah satu. \\
	\hline
	Menjalankan algoritma pilihan user & Solver dipilih berdasarkan input pengguna, lalu fungsi \texttt{solve} dipanggil terhadap \texttt{GameMap}. \\
	\hline
	Menghitung total cost & Setiap \texttt{MoveResult} menyimpan cost satu gerakan. Solver menjumlahkannya ke dalam \texttt{gCost}. \\
	\hline
	Menghitung jumlah iterasi & Setiap state yang diambil dari priority queue akan menambah counter iterasi. \\
	\hline
	Menghitung waktu eksekusi & Program menggunakan \texttt{std::chrono::high\_resolution\_clock} untuk menghitung waktu pencarian dalam milidetik. \\
	\hline
	Menampilkan solusi & Jika solusi ditemukan, \texttt{SearchResult} menyimpan path solusi, total cost, waktu eksekusi, dan jumlah iterasi. \\
	\hline
	Menyimpan solusi ke file & Jika pengguna memilih untuk menyimpan solusi, program menulis path solusi, cost, iterasi, waktu, dan visualisasi langkah ke folder \texttt{saves/}. \\
	\hline
	Playback CLI & Program dapat menampilkan ulang langkah solusi berdasarkan urutan path yang ditemukan. \\
	\hline
	GUI & Jika mode GUI dijalankan, program memvisualisasikan papan, aktor, obstacle, lava, angka, goal, dan playback solusi secara grafis. \\
	\hline
\end{longtable}

\section{Interface Solver}

Seluruh algoritma pathfinding menggunakan interface yang sama, yaitu \texttt{ISolver}. 
Interface ini membuat setiap algoritma memiliki fungsi \texttt{solve} dengan bentuk pemanggilan yang seragam. 
Dengan demikian, program utama dapat mengganti algoritma tanpa perlu mengubah alur besar pencarian.

\begin{tcolorbox}[
	colback=green!5,
	colframe=green!50!black,
	title=\textbf{Potongan Source Code Asli: \texttt{ISolver.hpp}},
	arc=2mm,
	boxrule=0.8pt
]
\small
\begin{verbatim}
#pragma once

#include "model/SearchResult.hpp"

class GameMap;

class ISolver {
public:
    virtual ~ISolver() = default;
    virtual SearchResult solve(const GameMap& map) const = 0;
};
\end{verbatim}
\end{tcolorbox}

\section{Implementasi Uniform Cost Search}

Class \texttt{UCS} mengimplementasikan algoritma Uniform Cost Search. 
Algoritma ini menggunakan \texttt{priority\_queue} dengan prioritas utama pada nilai \texttt{gCost}. 
Jika terdapat dua state dengan cost sama, panjang path digunakan sebagai pembanding tambahan.

\begin{tcolorbox}[
	colback=green!5,
	colframe=green!50!black,
	title=\textbf{Potongan Source Code Asli: Deklarasi \texttt{UCS.hpp}},
	arc=2mm,
	boxrule=0.8pt
]
\small
\begin{verbatim}
#pragma once

#include "solver/ISolver.hpp"

class GameMap;

class UCS : public ISolver {
public:
    UCS() = default;
    SearchResult solve(const GameMap& map) const override;
};
\end{verbatim}
\end{tcolorbox}

\begin{tcolorbox}[
	colback=green!5,
	colframe=green!50!black,
	title=\textbf{Potongan Source Code Asli: Prioritas dan Struktur Data UCS},
	arc=2mm,
	boxrule=0.8pt
]
\small
\begin{verbatim}
auto cmp = [](const State& a, const State& b) {
    if (a.getGCost() != b.getGCost()) return a.getGCost() > b.getGCost();
    return a.getPath().length() > b.getPath().length();
};
std::priority_queue<State, std::vector<State>, decltype(cmp)> openSet(cmp);
std::map<std::tuple<int, int, int>, int> bestCost;
\end{verbatim}
\end{tcolorbox}

\begin{table}[H]
	\centering
	\caption{Penjelasan implementasi UCS}
	\begin{tabular}{|p{4cm}|p{10cm}|}
		\hline
		\rowcolor{green!15}
		\textbf{Bagian Implementasi} & \textbf{Penjelasan} \\
		\hline
		\texttt{openSet} & Menyimpan state yang akan diproses, diurutkan berdasarkan cost aktual terkecil. \\
		\hline
		\texttt{bestCost} & Menyimpan cost terbaik untuk kombinasi posisi baris, posisi kolom, dan collected mask. \\
		\hline
		\texttt{iterations} & Menghitung banyaknya state yang diambil dari priority queue. \\
		\hline
		\texttt{MovementEngine::slide} & Membentuk successor state dengan simulasi sliding. \\
		\hline
		\texttt{SearchResult} & Mengembalikan informasi solusi, total cost, waktu eksekusi, dan jumlah iterasi. \\
		\hline
	\end{tabular}
\end{table}

\section{Implementasi Greedy Best-First Search}

Class \texttt{GBFS} menerima pointer ke \texttt{IHeuristic}. 
Heuristik digunakan untuk menentukan prioritas state. 
State dengan nilai heuristik lebih kecil akan diproses lebih dahulu. 
Jika nilai heuristik sama, program menggunakan \texttt{gCost} sebagai pembanding tambahan.

\begin{tcolorbox}[
	colback=orange!5,
	colframe=orange!70!black,
	title=\textbf{Potongan Source Code Asli: Deklarasi \texttt{GBFS.hpp}},
	arc=2mm,
	boxrule=0.8pt
]
\small
\begin{verbatim}
#pragma once

#include "solver/ISolver.hpp"

class IHeuristic;
class GameMap;

class GBFS : public ISolver {
public:
    explicit GBFS(const IHeuristic* heuristic);
    SearchResult solve(const GameMap& map) const override;

private:
    const IHeuristic* heuristic;
};
\end{verbatim}
\end{tcolorbox}

\begin{tcolorbox}[
	colback=orange!5,
	colframe=orange!70!black,
	title=\textbf{Potongan Source Code Asli: Prioritas GBFS},
	arc=2mm,
	boxrule=0.8pt
]
\small
\begin{verbatim}
auto cmp = [](const State& a, const State& b) {
    if (a.getHCost() != b.getHCost()) return a.getHCost() > b.getHCost();
    return a.getGCost() > b.getGCost();
};
std::priority_queue<State, std::vector<State>, decltype(cmp)> openSet(cmp);
std::map<std::tuple<int, int, int>, int> bestCost;
\end{verbatim}
\end{tcolorbox}

\begin{tcolorbox}[
	colback=orange!5,
	colframe=orange!70!black,
	title=\textbf{Potongan Source Code Asli: Perhitungan Heuristik pada GBFS},
	arc=2mm,
	boxrule=0.8pt
]
\small
\begin{verbatim}
State nextState(move.getNewPosition(),
                move.getNewCollectedMask(),
                move.getNewLastCollected(),
                current.getPath() + dir,
                newGCost,
                0,
                nullptr);
nextState.setHCost(heuristic->compute(nextState, map));
\end{verbatim}
\end{tcolorbox}

\begin{table}[H]
	\centering
	\caption{Penjelasan implementasi GBFS}
	\begin{tabular}{|p{4cm}|p{10cm}|}
		\hline
		\rowcolor{orange!20}
		\textbf{Bagian Implementasi} & \textbf{Penjelasan} \\
		\hline
		\texttt{heuristic} & Digunakan untuk menghitung estimasi kedekatan state terhadap target berikutnya. \\
		\hline
		\texttt{openSet} & Mengurutkan state berdasarkan nilai \texttt{hCost} terkecil. \\
		\hline
		\texttt{gCost} & Tetap disimpan untuk menghitung total cost solusi, tetapi bukan prioritas utama. \\
		\hline
		\texttt{bestCost} & Menghindari pemrosesan ulang state yang sudah dicapai dengan cost lebih kecil atau sama. \\
		\hline
	\end{tabular}
\end{table}

\section{Implementasi A*}

Class \texttt{AStar} juga menerima pointer ke \texttt{IHeuristic}. 
Perbedaan utamanya dengan GBFS adalah A* menggunakan nilai \texttt{fCost}, yaitu gabungan antara \texttt{gCost} dan \texttt{hCost}. 
Dengan demikian, A* mempertimbangkan cost aktual sekaligus estimasi menuju target.

\begin{tcolorbox}[
	colback=red!5,
	colframe=red!60!black,
	title=\textbf{Potongan Source Code Asli: Deklarasi \texttt{AStar.hpp}},
	arc=2mm,
	boxrule=0.8pt
]
\small
\begin{verbatim}
#pragma once

#include "solver/ISolver.hpp"

class IHeuristic;
class State;
class GameMap;

class AStar : public ISolver {
public:
    explicit AStar(const IHeuristic* heuristic);
    SearchResult solve(const GameMap& map) const override;

private:
    const IHeuristic* heuristic;
};
\end{verbatim}
\end{tcolorbox}

\begin{tcolorbox}[
	colback=red!5,
	colframe=red!60!black,
	title=\textbf{Potongan Source Code Asli: Prioritas A*},
	arc=2mm,
	boxrule=0.8pt
]
\small
\begin{verbatim}
auto cmp = [](const State& a, const State& b) {
    return a.getFCost() > b.getFCost();
};
std::priority_queue<State, std::vector<State>, decltype(cmp)> openSet(cmp);

std::map<std::tuple<int, int, int>, int> bestCost;
\end{verbatim}
\end{tcolorbox}

\begin{tcolorbox}[
	colback=red!5,
	colframe=red!60!black,
	title=\textbf{Potongan Source Code Asli: Ekspansi State pada A*},
	arc=2mm,
	boxrule=0.8pt
]
\small
\begin{verbatim}
MoveResult move = MovementEngine::slide(map, current, directions[i]);
if (!move.isValid() || move.isGameOver()) {
    continue;
}

std::string newPath = current.getPath() + directions[i];
int newGCost = current.getGCost() + move.getMoveCost();

State nextState(move.getNewPosition(),
               move.getNewCollectedMask(),
               move.getNewLastCollected(),
               newPath,
               newGCost,
               0,
               nullptr);
nextState.setHCost(heuristic->compute(nextState, map));
\end{verbatim}
\end{tcolorbox}

\begin{table}[H]
	\centering
	\caption{Penjelasan implementasi A*}
	\begin{tabular}{|p{4cm}|p{10cm}|}
		\hline
		\rowcolor{red!15}
		\textbf{Bagian Implementasi} & \textbf{Penjelasan} \\
		\hline
		\texttt{getFCost()} & Menghasilkan nilai prioritas A*, yaitu gabungan cost aktual dan heuristik. \\
		\hline
		\texttt{getGCost()} & Menyimpan total cost aktual dari start sampai state saat ini. \\
		\hline
		\texttt{setHCost()} & Mengatur nilai heuristik untuk state baru. \\
		\hline
		\texttt{bestCost} & Menyimpan cost terbaik untuk menghindari eksplorasi state yang tidak lebih baik. \\
		\hline
		\texttt{MovementEngine::slide} & Menentukan successor yang valid berdasarkan aturan sliding puzzle. \\
		\hline
	\end{tabular}
\end{table}

\section{Implementasi Weighted A*}

Class \texttt{WeightedAStar} merupakan implementasi algoritma tambahan sebagai spesifikasi bonus. 
Algoritma ini merupakan pengembangan dari A* dengan memberikan bobot lebih besar pada nilai heuristik. 
Pada implementasi program, Weighted A* menggunakan konstanta \texttt{WEIGHT = 2}, sehingga fungsi evaluasi yang digunakan adalah $g(n) + 2h(n)$. 
Dengan strategi ini, pencarian menjadi lebih agresif mengikuti arah target yang diperkirakan oleh heuristik.

\begin{tcolorbox}[
	colback=purple!5,
	colframe=purple!60!black,
	title=\textbf{Potongan Source Code Asli: Deklarasi \texttt{WeightedAStar.hpp}},
	arc=2mm,
	boxrule=0.8pt
]
\small
\begin{verbatim}
#pragma once

#include "solver/ISolver.hpp"

class IHeuristic;
class State;
class GameMap;

class WeightedAStar : public ISolver {
public:
    explicit WeightedAStar(const IHeuristic* heuristic);
    SearchResult solve(const GameMap& map) const override;

private:
    const IHeuristic* heuristic;
};
\end{verbatim}
\end{tcolorbox}

\begin{tcolorbox}[
	colback=purple!5,
	colframe=purple!60!black,
	title=\textbf{Potongan Source Code Asli: Bobot dan Prioritas Weighted A*},
	arc=2mm,
	boxrule=0.8pt
]
\small
\begin{verbatim}
constexpr int WEIGHT = 2;
auto cmp = [WEIGHT](const State& a, const State& b) {
    int fA = a.getGCost() + WEIGHT * a.getHCost();
    int fB = b.getGCost() + WEIGHT * b.getHCost();
    return fA > fB;
};
std::priority_queue<State, std::vector<State>, decltype(cmp)> openSet(cmp);

std::map<std::tuple<int, int, int>, int> bestCost;
\end{verbatim}
\end{tcolorbox}

Pada bagian tersebut, \texttt{priority\_queue} berfungsi sebagai \textit{open set}. 
State dengan nilai evaluasi paling kecil akan diproses lebih dahulu. 
Berbeda dengan A* biasa yang menggunakan $g(n)+h(n)$, Weighted A* menggunakan $g(n)+2h(n)$ sehingga nilai heuristik memiliki pengaruh lebih besar dalam pemilihan state.

\begin{table}[H]
	\centering
	\caption{Penjelasan struktur data Weighted A*}
	\begin{tabular}{|p{4cm}|p{10cm}|}
		\hline
		\rowcolor{purple!15}
		\textbf{Bagian Implementasi} & \textbf{Penjelasan} \\
		\hline
		\texttt{WEIGHT} & Konstanta bobot heuristik. Pada program digunakan nilai $2$. \\
		\hline
		\texttt{openSet} & Priority queue yang menyimpan state yang akan diproses berdasarkan nilai $g(n)+2h(n)$ terkecil. \\
		\hline
		\texttt{bestCost} & Map untuk menyimpan cost terbaik dari setiap kombinasi posisi aktor dan collected mask. \\
		\hline
		\texttt{heuristic} & Digunakan untuk menghitung nilai estimasi $h(n)$ pada state awal maupun state hasil ekspansi. \\
		\hline
	\end{tabular}
\end{table}

\begin{tcolorbox}[
	colback=purple!5,
	colframe=purple!60!black,
	title=\textbf{Potongan Source Code Asli: Inisialisasi State Awal Weighted A*},
	arc=2mm,
	boxrule=0.8pt
]
\small
\begin{verbatim}
Position startPos = map.getStartPosition();
State startState(startPos, 0, -1, "", 0, 0, nullptr);
startState.setHCost(heuristic->compute(startState, map));
openSet.push(startState);

int iterations = 0;
\end{verbatim}
\end{tcolorbox}

State awal dibentuk dari posisi awal aktor. 
Nilai \texttt{collectedMask} diatur menjadi $0$ karena belum ada angka yang dikumpulkan, nilai angka terakhir diatur menjadi $-1$, path masih kosong, dan cost awal bernilai $0$. 
Setelah itu, program menghitung nilai heuristik state awal dan memasukkannya ke dalam \texttt{openSet}.

\begin{tcolorbox}[
	colback=purple!5,
	colframe=purple!60!black,
	title=\textbf{Potongan Source Code Asli: Ekspansi State pada Weighted A*},
	arc=2mm,
	boxrule=0.8pt
]
\small
\begin{verbatim}
const char directions[4] = {'U', 'D', 'L', 'R'};
for (int i = 0; i < 4; ++i) {
    MoveResult move = MovementEngine::slide(map, current, directions[i]);
    if (!move.isValid() || move.isGameOver()) {
        continue;
    }

    std::string newPath = current.getPath() + directions[i];
    int newGCost = current.getGCost() + move.getMoveCost();

    State nextState(move.getNewPosition(),
                   move.getNewCollectedMask(),
                   move.getNewLastCollected(),
                   newPath,
                   newGCost,
                   0,
                   nullptr);
    nextState.setHCost(heuristic->compute(nextState, map));
\end{verbatim}
\end{tcolorbox}

Weighted A* mengekspansi state ke empat arah, yaitu \texttt{U}, \texttt{D}, \texttt{L}, dan \texttt{R}. 
Setiap arah tidak langsung dianggap valid, melainkan disimulasikan terlebih dahulu menggunakan \texttt{MovementEngine::slide}. 
Jika hasil gerakan tidak valid atau menyebabkan \textit{game over}, state tersebut langsung diabaikan. 
Jika gerakan valid, program membentuk path baru, menghitung \texttt{newGCost}, membuat state baru, lalu menghitung ulang nilai heuristiknya.

\begin{tcolorbox}[
	colback=purple!5,
	colframe=purple!60!black,
	title=\textbf{Potongan Source Code Asli: Pruning State Weighted A*},
	arc=2mm,
	boxrule=0.8pt
]
\small
\begin{verbatim}
auto nextKey = std::make_tuple(nextState.getPosition().getRow(),
                               nextState.getPosition().getCol(),
                               nextState.getCollectedMask());
auto nextIt = bestCost.find(nextKey);
if (nextIt != bestCost.end() && nextIt->second <= newGCost) {
    continue;
}

bestCost[nextKey] = newGCost;
openSet.push(nextState);
\end{verbatim}
\end{tcolorbox}

Bagian tersebut digunakan untuk menghindari eksplorasi state yang tidak lebih baik. 
Jika sudah ada state dengan posisi dan \texttt{collectedMask} yang sama serta cost lebih kecil atau sama, maka state baru tidak dimasukkan ke \texttt{openSet}. 
Sebaliknya, jika state baru memiliki cost yang lebih baik, \texttt{bestCost} diperbarui dan state tersebut dimasukkan ke priority queue.

\begin{table}[H]
	\centering
	\caption{Penjelasan implementasi Weighted A*}
	\begin{tabular}{|p{4cm}|p{10cm}|}
		\hline
		\rowcolor{purple!15}
		\textbf{Bagian Implementasi} & \textbf{Penjelasan} \\
		\hline
		\texttt{WEIGHT = 2} & Membuat heuristik memiliki pengaruh dua kali lebih besar dibanding A* biasa. \\
		\hline
		\texttt{getGCost()} & Menyimpan total cost aktual dari start sampai state saat ini. \\
		\hline
		\texttt{getHCost()} & Menyimpan nilai estimasi dari state saat ini menuju target berikutnya. \\
		\hline
		\texttt{gCost + WEIGHT * hCost} & Nilai evaluasi yang digunakan untuk mengurutkan state di priority queue. \\
		\hline
		\texttt{MovementEngine::slide} & Mensimulasikan gerakan sliding dan menghasilkan \texttt{MoveResult}. \\
		\hline
		\texttt{bestCost} & Menyimpan cost terbaik untuk menghindari pemrosesan ulang state yang lebih buruk. \\
		\hline
		\texttt{SearchResult} & Mengembalikan hasil pencarian, seperti status solusi, path solusi, total cost, waktu eksekusi, dan jumlah iterasi. \\
		\hline
	\end{tabular}
\end{table}

\begin{tcolorbox}[
	colback=purple!5,
	colframe=purple!60!black,
	title=\textbf{Karakteristik Implementasi Weighted A*},
	arc=2mm,
	boxrule=0.8pt
]
Weighted A* pada program ini memiliki struktur implementasi yang mirip dengan A*, tetapi berbeda pada fungsi evaluasi yang digunakan di priority queue. 
Dengan nilai $f(n)=g(n)+2h(n)$, algoritma menjadi lebih fokus terhadap estimasi menuju target. 
Hal ini dapat membuat pencarian lebih cepat dan jumlah iterasi lebih sedikit, tetapi solusi yang ditemukan tidak selalu memiliki total cost paling minimum.
\end{tcolorbox}

\section{Mekanisme Goal dan Hasil Pencarian}

Ketiga solver menggunakan kondisi goal yang sama. 
State dianggap sebagai solusi apabila posisi aktor sama dengan posisi goal dan seluruh angka telah dikumpulkan secara berurutan.

\begin{tcolorbox}[
	colback=blue!5,
	colframe=blue!60!black,
	title=\textbf{Potongan Source Code Asli: Pemeriksaan Goal dan SearchResult},
	arc=2mm,
	boxrule=0.8pt
]
\small
\begin{verbatim}
bool isGoal = current.getPosition() == map.getGoalPosition();
bool allNumbersCollected = (maxNumber == -1) || 
                            (current.getLastCollected() == maxNumber);

if (isGoal && allNumbersCollected) {
    auto endTime = std::chrono::high_resolution_clock::now();
    double elapsed = std::chrono::duration<double, std::milli>
                     (endTime - startTime).count();

    SearchResult result;
    result.setSolutionFound(true);
    result.setSolutionPath(current.getPath());
    result.setTotalCost(current.getGCost());
    result.setExecutionTimeMs(elapsed);
    result.setIterationsCount(iterations);
    return result;
}
\end{verbatim}
\end{tcolorbox}

\begin{table}[H]
	\centering
	\caption{Informasi yang disimpan pada hasil pencarian}
	\begin{tabular}{|p{4cm}|p{10cm}|}
		\hline
		\rowcolor{blue!15}
		\textbf{Data} & \textbf{Keterangan} \\
		\hline
		\texttt{solutionFound} & Menandakan apakah solusi berhasil ditemukan. \\
		\hline
		\texttt{solutionPath} & Menyimpan urutan gerakan solusi, misalnya \texttt{RULUDR}. \\
		\hline
		\texttt{totalCost} & Menyimpan total cost dari solusi. \\
		\hline
		\texttt{executionTimeMs} & Menyimpan waktu eksekusi algoritma dalam milidetik. \\
		\hline
		\texttt{iterationsCount} & Menyimpan banyaknya iterasi atau state yang ditinjau. \\
		\hline
		\texttt{message} & Menyimpan pesan jika solusi tidak ditemukan. \\
		\hline
	\end{tabular}
\end{table}

\section{Bukti Program Dapat Dijalankan}

Bagian ini menjelaskan perintah dasar untuk mengompilasi dan menjalankan program. 
Perintah dapat disesuaikan dengan target yang tersedia pada \texttt{Makefile}.

\begin{table}[H]
	\centering
	\caption{Perintah kompilasi dan eksekusi program}
	\begin{tabular}{|p{4cm}|p{10cm}|}
		\hline
		\rowcolor{purple!15}
		\textbf{Kebutuhan} & \textbf{Perintah} \\
		\hline
		Kompilasi program & \texttt{make} atau \texttt{make build} \\
		\hline
		Menjalankan CLI & \texttt{make run} atau \texttt{./build/ice\_solver} \\
		\hline
		Menjalankan GUI & \texttt{make gui}, jika Raylib tersedia dan target GUI diaktifkan. \\
		\hline
		Membersihkan hasil build & \texttt{make clean} \\
		\hline
	\end{tabular}
\end{table}

\subsection{Contoh Input File}

\begin{tcolorbox}[
	colback=gray!5,
	colframe=black!60,
	title=\textbf{Contoh File Input \texttt{.txt}},
	arc=2mm,
	boxrule=0.8pt
]
\small
\begin{verbatim}
7 7
XXXXXXX
X0****X
X**X**X
X****OX
X***1LX
XZ**X*X
XXXXXXX
999 999 999 999 999 999 999
999 3 5 2 8 1 999
999 7 4 999 6 9 999
999 2 8 3 5 4 999
999 6 1 7 2 999 999
999 9 3 4 999 8 999
999 999 999 999 999 999 999
\end{verbatim}
\end{tcolorbox}

\subsection{Contoh Output Singkat}

\begin{tcolorbox}[
	colback=gray!5,
	colframe=black!60,
	title=\textbf{Contoh Output Program},
	arc=2mm,
	boxrule=0.8pt
]
\small
\begin{verbatim}
>> Masukan file input:
test/input1.txt

>> Algoritma apa yang anda pilih? (UCS/GBFS/A*)
A*

>> Heuristic apa yang anda pilih? (H1/H2/H3)
H2

Solusi Yang Ditemukan : RULUDRUR
Cost dari Solusi : 87
Waktu eksekusi: 10 ms
Banyak iterasi yang dilakukan: 9999 iterasi

>> Apakah Anda ingin menyimpan solusi? (Ya/Tidak):
Ya

Solusi disimpan pada saves/solusi.txt
\end{verbatim}
\end{tcolorbox}
